\chapter*{Lời mở đầu}
\addcontentsline{toc}{chapter}{Lời mở đầu}
\label{ch:loi-mo-dau}

\section*{1. Lý do chọn đề tài}

Quy chế đào tạo là văn bản mà sinh viên tra cứu thường xuyên trong suốt quá trình học: thời gian tối đa
để hoàn thành khóa học, cách tính điểm học phần, điều kiện xét tốt nghiệp, thủ tục nghỉ học tạm thời hay
chuyển ngành. Tại Trường Đại học Sài Gòn (SGU), các quy định này nằm trong nhiều văn bản dài, phần lớn là
PDF bản scan, và có văn bản sửa đổi chỉ áp dụng cho một số khóa. Tìm theo từ khóa chỉ khớp khi người hỏi
dùng đúng thuật ngữ của văn bản (``cấm thi'' không tìm được ``không đủ điều kiện dự thi''); hỏi thẳng một
mô hình ngôn ngữ lớn (LLM) thì câu trả lời trôi chảy nhưng không có căn cứ và có thể bịa.

\textbf{LlamaIndex} được chọn vì nó chuyên cho đúng bài toán này --- hỏi đáp trên dữ liệu riêng --- và có
sẵn các khâu nạp dữ liệu, cắt Node, lập chỉ mục, truy hồi, tổng hợp câu trả lời~\cite{llamaindex2025docs};
hệ sinh thái của nó tích hợp sẵn Ollama (chạy LLM cục bộ, dữ liệu không ra dịch vụ bên ngoài) và mô hình
embedding tiếng Việt trên HuggingFace.

Định hướng nghề nghiệp của từng thành viên: \chualam{bổ sung ở bản báo cáo cuối}.

\section*{2. Mục tiêu báo cáo}

\paragraph{Mục tiêu kiến thức (cam kết).} Sáu mục lõi Tầng 1A:
\begin{enumerate}
  \item Ingestion: Document và metadata.
  \item Node/chunk, giải thích chunk size và overlap.
  \item Embedding và vector index.
  \item Retriever top-$k$ / similarity.
  \item Response synthesis có truy vết nguồn.
  \item Đánh giá thực nghiệm ở khâu truy hồi.
\end{enumerate}

\paragraph{Mục tiêu kỹ năng (kế hoạch).} Tầng 1B: custom node parser theo Điều, persistent index có
fingerprint, trích dẫn nguồn, từ chối câu hỏi ngoài corpus, nhiều nguồn (6 văn bản). Tầng 3: hybrid
BM25 + vector, xếp hạng lại bằng cross-encoder, đánh giá có khoảng tin cậy và kiểm định thống kê, model
server cho embedding/reranker.

\section*{3. Phạm vi nghiên cứu}

\begin{itemize}
  \item \textbf{Corpus}: 6 văn bản SGU đang áp dụng --- Quy chế đào tạo trình độ đại học 2021, QĐ 3183/2025
        sửa đổi Quy chế, QĐ 810 (học cùng lúc hai chương trình), QĐ 2375/2022 (chuyển chương trình, ngành),
        quy định khối lượng đào tạo, quy định số tín chỉ tối thiểu kỹ sư CNTT --- tổng 54 Node.
  \item \textbf{Mô hình}: embedding \texttt{AITeamVN/Vietnamese\_Embedding} chạy local (GPU 4\,GB ở fp16 hoặc
        CPU); LLM \texttt{qwen3:8b} qua Ollama trên máy GPU thuê hoặc Google Colab.
  \item \textbf{Không làm}: tư vấn pháp lý, agent, fine-tune LLM. Giao diện web chỉ ở mức chạy thử (Gradio).
\end{itemize}

\paragraph{Thay đổi phạm vi so với kế hoạch ban đầu.} Kế hoạch ban đầu là một văn bản (Quy chế 2021, 17
trang scan) với LLM \texttt{qwen3:4b} chạy hoàn toàn trên máy cá nhân. Corpus được mở rộng lên 6 văn bản
để trả lời được nội dung sửa đổi; LLM chuyển sang chạy từ xa vì GPU 4\,GB không chứa nổi
\texttt{qwen3:8b} cùng embedding. Lý do chi tiết ở Chương~\ref{ch:danh-gia}.

\section*{4. Kế hoạch và tiến độ 12 tuần}

\begin{table}[H]
  \centering
  \small
  \begin{tabularx}{\textwidth}{l >{\hsize=0.9\hsize}X >{\hsize=1.3\hsize}X >{\hsize=0.8\hsize}X}
    \toprule
    \textbf{Tuần} & \textbf{Nội dung chính} & \textbf{Việc cụ thể của nhóm} & \textbf{Trạng thái} \\
    \midrule
    1 & Lập nhóm, chọn công nghệ & Đăng ký LlamaIndex, hồ sơ RAG & Xong \\
    2 & Môi trường, Hello World, repo & venv, Ollama, \texttt{hello\_world.py}, \texttt{check\_model\_memory.py} & Xong \\
    3--4 & Bắt buộc lõi, nháp Ch.~3--4 & OCR, Node theo Điều, index; bộ câu hỏi 01/02 & Xong; thay bằng Markdown hai tầng \\
    5 & Chốt bài toán, thiết kế luồng & 6 văn bản, sơ đồ pipeline; Colab; chạy lại 56 câu & Xong \\
    6--7 & Chọn thêm, thực nghiệm & Retriever, query engine, Recall@$k$ (xong sớm); thí nghiệm chunking & Đang làm \\
    8--9 & Tích hợp, Tầng 3 & Trích nguồn, rerank, hybrid (xong sớm); từ chối ngoài corpus & Một phần \\
    10 & Kiểm thử, đo trước--sau & 26 kiểm thử + CI (xong sớm); chấm câu trả lời tự động & Một phần \\
    11 & Hands-on lab, chấm chéo & Viết lab & Chưa \\
    12 & Hoàn thiện, niêm phong & Báo cáo, slide, video; nộp $\ge$ 24 giờ trước bảo vệ & Chưa \\
    \bottomrule
  \end{tabularx}
  \caption{Kế hoạch 12 tuần và trạng thái tại tuần 5. Người thực hiện từng tuần ghi ở Phụ lục~\ref{app:nhat-ky}.}
\end{table}

\section*{5. Cấu trúc báo cáo}

Chương 1--2 giới thiệu LlamaIndex, kiến trúc, so sánh dựa trên trải nghiệm của nhóm
và phạm vi nên/không nên dùng. Chương 3--8 trình bày môi trường, sáu mục lõi Tầng 1A
với mã nguồn và kết quả chạy thật, rồi các kỹ thuật chọn thêm (node parser, persistent index, trích
dẫn/từ chối) và nâng cao (hybrid, rerank). Chương 9--10 là phân tích thiết kế, triển
khai, minh chứng và Bảng phủ kỹ thuật. Chương 11--12 đánh giá lại so với mục tiêu, phân
tích lỗi, AI Disclosure và hướng phát triển. Cuối báo cáo là tài liệu tham khảo và các phụ lục: nhật ký
tuần, bảng tự kiểm ba tầng và thực nghiệm học tập.

Mọi số liệu đều lấy từ việc chạy thật các script trong repository; phần chưa làm được đánh dấu bằng
khung ``Chưa triển khai'' hoặc ghi chú ``Chưa làm'' kèm thời điểm dự kiến bổ sung.
